\documentclass[11pt,a4paper]{article}
\usepackage[utf8]{inputenc}
\usepackage[T1]{fontenc}
\usepackage[ngerman]{babel}
\usepackage{lmodern}
\usepackage{microtype}
\usepackage[a4paper,margin=22mm,top=23mm,bottom=24mm]{geometry}
\usepackage{amsmath,amssymb,mathtools}
\usepackage{booktabs,longtable,tabularx,array}
\usepackage{xcolor}
\usepackage[most]{tcolorbox}
\usepackage{enumitem}
\usepackage{hyperref}
\usepackage{fancyhdr}
\usepackage{lastpage}
\usepackage{titlesec}
\usepackage{parskip}
\usepackage{ragged2e}

\definecolor{navy}{HTML}{17365D}
\definecolor{blue}{HTML}{2F5597}
\definecolor{green}{HTML}{2E7D32}
\definecolor{amber}{HTML}{A66A00}
\definecolor{red}{HTML}{B3261E}
\definecolor{purple}{HTML}{6A3D9A}
\definecolor{graydark}{HTML}{4A4A4A}
\definecolor{graylight}{HTML}{F2F4F7}
\definecolor{bluelight}{HTML}{EAF1FB}
\definecolor{greenlight}{HTML}{EAF5EA}
\definecolor{amberlight}{HTML}{FFF5DD}
\definecolor{redlight}{HTML}{FCEBEC}
\definecolor{purplelight}{HTML}{F3ECF9}

\hypersetup{colorlinks=true,linkcolor=blue,urlcolor=blue,pdftitle={Lighthouse-JAX Master Report v0.1},pdfauthor={Lighthouse-JAX MASTER}}
\setlist[itemize]{leftmargin=5mm,itemsep=2pt,topsep=3pt}
\setlist[enumerate]{leftmargin=6mm,itemsep=2pt,topsep=3pt}
\renewcommand{\arraystretch}{1.18}
\setlength{\parindent}{0pt}
\setlength{\emergencystretch}{2em}

\titleformat{\section}{\Large\bfseries\color{navy}}{\thesection}{0.6em}{}
\titleformat{\subsection}{\large\bfseries\color{blue}}{\thesubsection}{0.6em}{}

\pagestyle{fancy}
\fancyhf{}
\lhead{\small Lighthouse-JAX -- MASTER Report v0.1}
\rhead{\small 04.09.2026}
\cfoot{\small Seite \thepage\ von \pageref{LastPage}}
\renewcommand{\headrulewidth}{0.4pt}

\newcommand{\status}[2]{\tcbox[on line,boxrule=0pt,arc=1.5pt,left=4pt,right=4pt,top=2pt,bottom=2pt,colback=#2!12,coltext=#2,fontupper=\bfseries\footnotesize]{#1}}
\newcommand{\STABLE}{\status{STABLE}{green}}
\newcommand{\OPEN}{\status{OPEN}{amber}}
\newcommand{\WAIT}{\status{WAIT}{graydark}}
\newcommand{\PROTECTED}{\status{PROTECTED}{purple}}
\newcommand{\FROZEN}{\status{FROZEN}{blue}}
\newcommand{\FAILED}{\status{FAILED}{red}}
\newcommand{\ASSUMPTION}{\status{ASSUMPTION}{amber}}
\newcommand{\CONJECTURE}{\status{CONJECTURE}{purple}}

\newtcolorbox{summarybox}{colback=bluelight,colframe=blue,boxrule=0.7pt,arc=2pt}
\newtcolorbox{warningbox}{colback=amberlight,colframe=amber,boxrule=0.7pt,arc=2pt}
\newtcolorbox{successbox}{colback=greenlight,colframe=green,boxrule=0.7pt,arc=2pt}
\newtcolorbox{stopbox}{colback=redlight,colframe=red,boxrule=0.9pt,arc=2pt}

\begin{document}

\begin{titlepage}
\thispagestyle{empty}
\vspace*{18mm}
{\Huge\bfseries\color{navy} Lighthouse-JAX\par}
\vspace{4mm}
{\LARGE MASTER Project Report\par}
\vspace{3mm}
{\Large Version 0.1\par}
\vspace{14mm}
\begin{tcolorbox}[colback=graylight,colframe=navy,boxrule=0.8pt,arc=3pt]
\textbf{Projekt:} Modern mathematical and computational reconstruction of Hermann Haken's Lighthouse model\\[2mm]
\textbf{Repository:} \texttt{twkroll/haken-lighthouse-jax}\\[2mm]
\textbf{Berichtsdatum:} 04.09.2026\\[2mm]
\textbf{Kanonischer Projektstatus:} v0.2\\[2mm]
\textbf{Aktiver wissenschaftlicher Gate:} CORE Mathematical Scope Gate 0.1
\end{tcolorbox}
\vfill
\begin{summarybox}
\textbf{Governance-Hinweis.} Dieser Bericht ist ein Snapshot des kanonischen Git-Zustands. Er erzeugt keine neue wissenschaftliche Theorie. Nicht eingefrorene Inhalte werden als \OPEN, \WAIT\ oder \PROTECTED\ gekennzeichnet.
\end{summarybox}
\vspace{10mm}
{\small Erstellt durch 00 -- MASTER im Rahmen des Projekt-Governance-Protokolls.}
\end{titlepage}

\tableofcontents
\newpage

\section{Executive Summary}

\begin{summarybox}
Das Projekt befindet sich nach abgeschlossener Governance-Initialisierung vor seinem ersten wissenschaftlichen Gate. Die Governance-Regeln sind \STABLE. Der mathematische Kern des Lighthouse-Modells ist dagegen noch \OPEN. Deshalb sind JAX-Implementierung, Parameterstudien, Anwendungsbenchmarks, Differentiable-Inference-Experimente, Literatur-Novelty-Positionierung und Manuskriptarbeit derzeit nicht autorisiert.
\end{summarybox}

Die zentrale nächste Aufgabe ist eindeutig: Im Chat \texttt{10 -- CORE -- Haupttheorie / mathematischer Kern} muss exakt \texttt{GO} ausgeführt werden. Der dort bereits eingefrorene Auftrag verlangt ausschließlich die mathematische Scope-Klärung des Basismodells und endet mit \textbf{STOP -- RETURN TO MASTER}.

Der aktuelle Projektstand enthält noch keine wissenschaftlich eingefrorenen Resultate. Entsprechend gibt es noch keine STRONG-, WEAK-, NULL- oder FAIL-Ergebnisse, keine branch-abhängigen Claims und keinen Manuskript-Claim-Freeze. Der letzte stabile Rollback-Punkt ist \textbf{RB-001 Governance Initialization 0.1}.

\section{Zentrale Forschungsfrage}

Die kanonische Forschungsfrage lautet:

\begin{quote}
\textit{Can Hermann Haken's Lighthouse model be developed into a modern, scalable and differentiable framework for spiking-network dynamics that remains mathematically analyzable while supporting contemporary numerical simulation, inference and neuromorphic applications?}
\end{quote}

Das Projekt ist in vier geplante Ebenen gegliedert:
\begin{enumerate}
\item historischer mathematischer Kern,
\item präzise mathematische Rekonstruktion,
\item moderne Erweiterungen wie Graphen, Kontinuumsgrenzen, Delays, Rauschen, Plastizität und differenzierbare Approximationen,
\item reproduzierbare, vektorisierte und accelerator-fähige JAX-Implementierung.
\end{enumerate}

Diese Ebenen sind derzeit eine Projektarchitektur, keine Aussage darüber, welche Erweiterungen später wissenschaftlich erfolgreich sein werden.

\section{Aktueller Projektstatus}

\begin{tabularx}{\textwidth}{>{\bfseries}p{0.19\textwidth} p{0.23\textwidth} X}
\toprule
Workstream & Status & Aktuelle Rolle \\
\midrule
00 MASTER & \FROZEN\ / \WAIT & Governance-Aufsicht; wartet auf CORE-Ergebnis \\
10 CORE & \status{READY}{green} & Erster mathematischer Scope-Gate \\
50 APP-1 Neuroscience & \PROTECTED\ / \WAIT & Reservierter Anwendungsbranch \\
60 APP-2 Neuromorphic & \PROTECTED\ / \WAIT & Reservierter Anwendungsbranch \\
70 APP-3 Diff. Inference & \PROTECTED\ / \WAIT & Reservierter Anwendungsbranch \\
80 LIT & \WAIT & Literatur-Audit erst nach MASTER-Autorisierung \\
90 MANUSCRIPT & \WAIT & Keine Claims/Figuren vor frozen Resultaten \\
\bottomrule
\end{tabularx}

\subsection{Aktiver Blocker}

Der kanonische mathematische Scope des Basismodells wurde noch nicht etabliert und eingefroren. In \texttt{research/core/STATUS.md} steht CORE weiterhin auf \texttt{READY}; eine Ergebnisdatei \texttt{research/core/mathematical\_scope\_gate\_0\_1.md} existiert noch nicht.

\subsection{Freeze-Check}

\begin{successbox}
\textbf{Freeze-Check: OK.} Bis zum aktuellen Snapshot wurden keine Effect-Inspections, Parameter-Retunings, Objective-Wechsel, Anwendungs-Executions oder post-hoc Auswahlen dokumentiert. Der README-only Commit \texttt{948dedbc5294fbe864b940060ee6b2053020347f} wurde als nicht-wissenschaftliche Execution klassifiziert.
\end{successbox}

\section{Theoretische Grundlagen und mathematischer Freeze-Status}

\subsection{Was bereits stabil ist}

Stabil ist derzeit nicht eine konkrete Modellgleichung, sondern das Verfahren, wie diese Gleichung ausgewählt und eingefroren werden muss. Der wissenschaftliche Workflow ist formal:

\[
\text{Gate} \longrightarrow \text{Freeze} \longrightarrow \text{Execution} \longrightarrow \text{Result Freeze}.
\]

Für das aktuelle Projekt gilt damit:

\[
\text{CORE Mathematical Scope Gate 0.1}
\longrightarrow
\text{CORE Mathematical Freeze 0.1}
\longrightarrow
\text{spätere numerische Spezifikation}.
\]

\subsection{Was bewusst noch nicht kanonisch ist}

\begin{warningbox}
\textbf{Mathematische Modellgleichungen: \OPEN.} Das Projekt hat noch keine kanonische Gleichungsmenge, keine finalen Zustandsvariablen, keine eingefrorene Schwellenfunktion, keine eingefrorene Synapsenform und keine eingefrorene Delay-Konvention. Diese Inhalte werden absichtlich nicht aus früheren Chat-Antworten oder aus dem README rekonstruiert, weil Git die Single Source of Truth ist und der CORE-Gate noch nicht ausgeführt wurde.
\end{warningbox}

Formal ist die Menge der bereits wissenschaftlich eingefrorenen CORE-Aussagen zum Snapshot leer:

\[
\mathcal{R}_{\mathrm{CORE}}^{\mathrm{frozen}} = \varnothing.
\]

Dies bedeutet nicht, dass es keine bekannte Lighthouse-Theorie gibt; es bedeutet nur, dass innerhalb dieses Projekts noch keine solche Definition als kanonischer Freeze akzeptiert wurde.

\subsection{Scope des nächsten CORE-Gates}

Der autorisierte Gate muss unter anderem:
\begin{itemize}
\item primäre Haken-Formulierungen von späteren Reformulierungen unterscheiden,
\item kanonische Notation für Zustand, Phasenentwicklung, Spike-Ereignisse, synaptische Aktivierung, Kopplung, Schwellen und Delays definieren,
\item das minimale Basismodell für spätere numerische Reproduktion festlegen,
\item zulässige Modellvarianten registrieren, ohne nach beobachteter Effektstärke zu selektieren,
\item etablierte Aussagen, Projektannahmen, Conjectures und Open Questions trennen,
\item analytische Validierungsziele wie Synchronisation, Phase Locking, Delay-Effekte, Wellen, Bumps, Stabilität und Kontinuumsgrenzen nur quellenbasiert dokumentieren.
\end{itemize}

\section{CORE-Ergebnisse}

\begin{warningbox}
\textbf{Status: OPEN / NONE YET.} Es existiert noch kein eingefrorenes CORE-Ergebnis. Daher werden in dieser Version weder Theoreme noch Propositionen, Modellparameter oder numerische Referenzresultate behauptet.
\end{warningbox}

Vorgesehene Ergebnisdatei des nächsten Gates: \texttt{research/core/mathematical\_scope\_gate\_0\_1.md}.

\section{Anwendungsbranches}

\subsection{APP-1 Computational Neuroscience}
Status: \PROTECTED\ / \WAIT. Abhängigkeiten: CORE Mathematical Freeze und späteres MASTER-Anwendungsgate. Es existieren noch keine Kandidaten, Parameter, Objectives oder Resultate.

\subsection{APP-2 Neuromorphic Computing}
Status: \PROTECTED\ / \WAIT. Die bloße Projektidee einer späteren accelerator-/hardware-nahen Implementierung ist kein Anwendungsresultat.

\subsection{APP-3 Differentiable Inference / Temporal Learning}
Status: \PROTECTED\ / \WAIT. Differenzierbare Approximationen oder Lernexperimente sind derzeit nicht autorisiert; es gibt keine gefrorenen Losses, Datensätze oder Erfolgsmetriken.

\section{Literaturpositionierung}

Der Literaturbranch steht auf \WAIT. Ein gezielter Novelty-Audit ist erst nach MASTER-Autorisierung vorgesehen. Entsprechend sind Novelty-Stufen N0--N4 noch nicht vergeben.

\subsection{Referenzkandidaten aus dem Repository-README}
Die folgenden Quellen sind im README als Core References verzeichnet, wurden in diesem PDF-Snapshot aber nicht als vollständiger Literatur-Audit neu bewertet:
\begin{enumerate}
\item H. Haken, \textit{Quasi-discrete dynamics of a neural net: The Lighthouse model}, Discrete Dynamics in Nature and Society 4 (2000), 187--200. DOI: 10.1155/S1026022600000182.
\item H. Haken, \textit{Phase Locking in the Lighthouse Model of a Neural Net with Several Delay Times}, Progress of Theoretical Physics Supplement 139 (2000), 96--111. DOI: 10.1143/PTPS.139.96.
\item H. Haken, \textit{Brain Dynamics: Synchronization and Activity Patterns in Pulse-Coupled Neural Nets with Delays and Noise}, Springer (2002).
\item S. Coombes, \textit{Revisiting the Haken Lighthouse model}, European Physical Journal Special Topics, version of record 2025, volume 235 (2026), 4571--4593. DOI: 10.1140/epjs/s11734-025-01841-3.
\item S. Coombes, R. Thul, S. Ruschel, R. Nicks, \textit{Adaptive conduction delays and phase locking in spiking Haken Lighthouse networks}, arXiv:2606.21508 (2026).
\end{enumerate}

\begin{warningbox}
Kein SAME-/CLOSE-/RELATED-Rating ist derzeit eingefroren. Aus der Existenz oder Abwesenheit einzelner Treffer wird keine Novelty-Aussage abgeleitet.
\end{warningbox}

\section{Stabile Resultate, Annahmen und offene Fragen}

\subsection{STABLE}
\begin{itemize}
\item Governance-Prozess 0.1: \STABLE.
\item Git-Repository \texttt{twkroll/haken-lighthouse-jax} als Single Source of Truth: \STABLE.
\item Workstream-Struktur 00/10/50/60/70/80/90: \STABLE.
\item Anwendungsbranches bleiben bis MASTER-Gate \PROTECTED.
\end{itemize}

\subsection{ASSUMPTION / Projektprämissen}
\begin{itemize}
\item Die Modernisierung soll mathematische Analysierbarkeit nicht ohne explizite Entscheidung zugunsten bloßer Rechenleistung aufgeben.
\item JAX ist eine geplante Implementierungsrichtung, aber noch keine eingefrorene numerische Spezifikation.
\end{itemize}

\subsection{OPEN}
\begin{itemize}
\item kanonische Lighthouse-Modellgleichungen und Notation,
\item zulässige Varianten und Exklusionen,
\item numerisches Integrations-/Event-Konzept,
\item Delay- und Synapsenrepräsentation,
\item Differentiationsstrategie,
\item konkrete Anwendungs-Candidates,
\item Novelty-Klassifikation und Manuskriptclaim.
\end{itemize}

\subsection{WEAK / NULL / FAIL}
Keine eingefrorenen Resultate vorhanden. Insbesondere werden negative oder schwache Resultate nicht vorweggenommen und später nicht aus dem Bericht ausgeblendet.

\section{Branch-independent vs. branch-dependent}

\textbf{Branch-independent:}
\begin{itemize}
\item Governance-Workflow und Anti-Cherry-Picking-Regeln,
\item Git als kanonischer Zustands-/Prompt-/Freeze-Speicher,
\item Reihenfolge des ersten CORE-Gates vor Implementierung und Anwendung.
\end{itemize}

\textbf{Branch-dependent:} derzeit keine wissenschaftlichen Resultate.

\section{Decision \& Branch Log}

\begin{longtable}{p{0.12\textwidth} p{0.17\textwidth} p{0.61\textwidth}}
\toprule
ID & Status & Entscheidung \\
\midrule
DEC-001 & \STABLE & Governance-Prompt als Projektverfahren übernommen. \\
DEC-002 & \STABLE & Git-Repository ist Single Source of Truth. \\
DEC-003 & \STABLE & Initiale Workstream-Struktur autorisiert. \\
DEC-004 & \FROZEN & Application-Workstreams nur PROTECTED / WAIT bis explizites MASTER-Gate. \\
DEC-005 & \status{ACTIVE}{blue} & Erster wissenschaftlicher Auftrag ist CORE Mathematical Scope Gate 0.1; keine Implementierung oder Anwendung davor. \\
DEC-006 & \STABLE & MASTER Status Audit 0.2 bestätigt: CORE-Gate noch nicht ausgeführt; README-Commit ist kein Freeze-Verstoß. \\
DEC-007 & \STABLE & MASTER PDF Snapshot v0.1 erzeugt; wissenschaftlicher Freeze-Zustand bleibt unverändert. \\
\bottomrule
\end{longtable}

\section{Freeze-Status und Rollback-Punkte}

\begin{tabularx}{\textwidth}{X p{0.25\textwidth}}
\toprule
Objekt & Zustand \\
\midrule
Governance rules 0.1 & \STABLE \\
Scientific model definition & \OPEN \\
Numerical specification & \OPEN \\
Application candidates & \status{NOT AUTHORIZED}{red} \\
Claims / novelty & \OPEN \\
CORE Mathematical Scope Gate 0.1 & \status{READY}{green} \\
\bottomrule
\end{tabularx}

\subsection{Rollback}
\textbf{RB-001 Governance Initialization 0.1 -- STABLE.}

Weitere Rollback-Punkte werden erst nach tatsächlichen wissenschaftlichen Freezes ergänzt; sie dürfen RB-001 nicht rückwirkend überschreiben.

\section{Manuskriptstatus}

Status: \WAIT. Kein Claim Freeze, kein Draft, keine Figur und kein Supplement sind autorisiert. Der Manuskriptchat darf später ausschließlich auf bereits freigegebenen/frozen Resultaten aufbauen, sofern MASTER nicht ausdrücklich neue Wissenschaft autorisiert.

\section{Roadmap, Blocker und Abhängigkeiten}

\begin{enumerate}
\item \textbf{Jetzt:} CORE Mathematical Scope Gate 0.1 ausführen.
\item Bei PASS: vorgeschlagenen CORE Mathematical Freeze 0.1 durch MASTER prüfen und gegebenenfalls autorisieren.
\item Danach: numerische Spezifikation als eigener Gate/Freeze, bevor Implementierungsergebnisse inspiziert werden.
\item Erst danach: kontrollierte JAX-Execution und Referenzvalidierung.
\item Applications bleiben PROTECTED, bis der mathematische und numerische Kern stabil genug ist.
\item Literatur-Novelty-Audit gezielt nach einem geeigneten Result Freeze oder nach expliziter MASTER-Entscheidung.
\item Manuskript erst aus frozen Claims und Ergebnissen.
\end{enumerate}

\textbf{Aktiver Blocker:} fehlender CORE Mathematical Scope Gate 0.1.

\textbf{Abhängigkeit:}
\[
\text{Applications / JAX / Literature / Manuscript}
\quad \not\prec \quad
\text{CORE Scope Gate 0.1}.
\]
Das Symbol \(\not\prec\) wird hier als Governance-Kurznotation verwendet: Diese Arbeitsstränge dürfen den CORE-Gate nicht zeitlich in der wissenschaftlichen Execution überholen.

\section{Nächster globaler Schritt}

\begin{stopbox}
\textbf{Nächster Schritt:} Wechsle in \texttt{10 -- CORE -- Haupttheorie / mathematischer Kern} und schreibe dort exakt:

\begin{center}
\Large\texttt{GO}
\end{center}

Nach Abschluss und Git-Update: zurück zu \texttt{00 -- MASTER} und exakt \texttt{Status?}.

\textbf{STOP -- AWAIT CORE GO}
\end{stopbox}

\section{Versionhistorie}

\begin{tabularx}{\textwidth}{p{0.13\textwidth} p{0.18\textwidth} X}
\toprule
Version & Datum & Änderung \\
\midrule
v0.1 & 04.09.2026 & Erster kanonischer MASTER PDF Snapshot. Grundlage: Projektstatus v0.2, Branch-STATUS-Dateien, Decision \& Branch Log und Repository-README. Keine Änderung des wissenschaftlichen Freeze-Zustands. \\
\bottomrule
\end{tabularx}

\section*{Kanonische Repository-Quellen für diesen Snapshot}
\addcontentsline{toc}{section}{Kanonische Repository-Quellen}
\begin{itemize}
\item \texttt{research/master/project\_status.md} (v0.2)
\item \texttt{research/master/STATUS.md}
\item \texttt{research/master/decision\_branch\_log.md}
\item \texttt{research/master/prompts/core\_mathematical\_scope\_gate\_0\_1.md}
\item \texttt{research/core/STATUS.md}
\item \texttt{research/app\_neuroscience/STATUS.md}
\item \texttt{research/app\_neuromorphic/STATUS.md}
\item \texttt{research/app\_diff\_inference/STATUS.md}
\item \texttt{research/literature/STATUS.md}
\item \texttt{research/manuscript/STATUS.md}
\item \texttt{README.md}
\end{itemize}

\end{document}
